% Gerado por scripts/etapa3_comparacao.py — requer \usepackage{booktabs}
\begin{table}[htbp]
\centering
\small
\caption{Combinação HW + XGBoost — quantis empíricos versus normais (sensibilidade: estoque arredondado para cima)}
\label{tab:comb_empirico_vs_normal_sensibilidade}
\begin{tabular}{lrrrrrrr}
\toprule
Nível (razão) & Custo empírico & Custo normal & Diferença [IC 95\%] & Sobra empírico & Sobra normal & Falta empírico & Falta normal \\
\midrule
80\% (1:4) & 2,722 & 2,838 & +4,3 [+3,8; +4,8] & 1,766 & 2,092 & 0,239 & 0,187 \\
90\% (1:9) & 3,823 & 3,824 & +0,0 [-0,3; +0,4] & 2,660 & 2,768 & 0,129 & 0,117 \\
95\% (1:19) & 4,962 & 4,895 & -1,4 [-1,6; -1,1] & 3,615 & 3,345 & 0,071 & 0,082 \\
\bottomrule
\end{tabular}
\par\smallskip
{\footnotesize Nota: teste de 90 dias (26/01 a 24/04/2016), 5.748 séries. Diferença: custo normal em relação ao empírico, em \% (positivo = normal mais caro), com intervalo de 95\% de confiança entre colchetes, item como unidade (1.437 itens). Sobra e falta em unidades por série-dia. custo por série-dia com custo de sobra = 1; razão 1:4, 1:9 e 1:19 (sobra:falta) corresponde aos níveis de serviço de 80\%, 90\% e 95\%. Estoque = quantil arredondado para cima.}
\end{table}
